% 从旧数学附录迁出的剩余、Lucas 与阶乘表说明；不新增算法。
\section{素数模平方根}\label{knowledge-quadratic-residue}\index{Cipolla}\index{quadratic residue}
奇素数 $p$、$n\not\equiv0$：有平方根当且仅当 $n^{(p-1)/2}\equiv1\pmod p$。
$p\equiv3\pmod4$ 时 $r=n^{(p+1)/4}$，另一根 $p-r$。
Cipolla：取 $a$ 使 $w=a^2-n$ 非平方，$\omega^2=w$，根为 $(a+\omega)^{(p+1)/2}$。
实现见第~\pageref{compact-mod_sqrt}~页；零和模 2 单独处理；合数模不适用。

\section{乘法阶与原根}\label{knowledge-primitive-root}\index{原根}\index{multiplicative order}
$\gcd(a,n)=1$ 时阶 $d\mid\varphi(n)$：从 $d=\varphi(n)$ 起，对每个素因子 $q$，$a^{d/q}\equiv1$ 就除去 $q$。
不可逆 $a$ 接口返回 0。
$a$ 为原根当且仅当 $\gcd(a,n)=1$ 且对 $\varphi(n)$ 每个素因子 $q$ 有 $a^{\varphi(n)/q}\not\equiv1$。
$n\ge2$ 时原根仅存在于 $2,4,p^k,2p^k$（$p$ 奇素数）。
$\operatorname{ord}(g^t)=\varphi(n)/\gcd(t,\varphi(n))$，全部原根为 $\gcd(t,\varphi(n))=1$ 的 $g^t$，
共 $\varphi(\varphi(n))$ 个，按大小输出需排序。实现见第~\pageref{compact-PrimitiveRoot}~页。
$\operatorname{ord}(a^k)=d/\gcd(d,k)$；两阶互素时乘积之阶为两阶之积。

\section{Carmichael：最小公共指数}\label{knowledge-carmichael}\index{Carmichael}\index{Korselt criterion}
$\lambda(n)$ 为所有可逆元阶的最小公倍数，$\lambda(1)=1$，$\lambda(n)\mid\varphi(n)$。
\[
\lambda(n)=\operatorname{lcm}_i\lambda(p_i^{e_i}),\qquad
\lambda(p^e)=
\begin{cases}
2^{e-2},&p=2,\ e\ge3,\\
p^{e-1}(p-1),&\text{其他情形}.
\end{cases}
\]
仅可逆底数可把指数对 $\lambda(n)$ 取模。
合数 $n$ 是 Carmichael 数当且仅当 $\lambda(n)\mid(n-1)$，等价于无平方因子且每个素因子 $p$ 满足 $(p-1)\mid(n-1)$（须排除素数）；
Fermat 检验不能代替素性判定。

\section{Lucas：按素数进制逐位计算}\label{knowledge-lucas}\index{Lucas}
$p$ 为素数，$n_i,k_i$ 为 $p$ 进制位：$\binom nk\equiv\prod_i\binom{n_i}{k_i}\pmod p$，某位 $k_i>n_i$ 则为零。
只预处理 $0!,\ldots,(p-1)!$ 及逆元；每位常数次乘法。
实现见第~\pageref{compact-Lucas}~页，参数为 $n,k$；自行求和时检查 64 位溢出。

\section{合数模数组合数与插值}\label{knowledge-exlucas}\index{exLucas}\index{Lagrange interpolation}
$q=p^e$，$v_p(n!)=\sum_{j\ge1}\lfloor n/p^j\rfloor$，$U_p(n)=n!/p^{v_p(n!)}$，
$t=v_p(n!)-v_p(k!)-v_p((n-k)!)$。$t\ge e$ 时为零，否则
\[
\binom nk\equiv U_p(n)\,U_p(k)^{-1}\,U_p(n-k)^{-1}\,p^t\pmod q,
\]
各素数幂结果 CRT 合并，见第~\pageref{compact-ExLucas}~页。

素数域中 $n$ 个两两不同的 $x_i$，$w_i=y_i/\prod_{j\ne i}(x_i-x_j)$：
\[
f(k)=\sum_{i=0}^{n-1}w_i\prod_{j<i}(k-x_j)\prod_{j>i}(k-x_j),
\]
查询点可等于样本点。$x_i=i$ 时分母为 $i!\,(n-1-i)!\,(-1)^{n-1-i}$，要求 $n\le p$。
见第~\pageref{compact-Lagrange}~页。$S_d(k)=\sum_{i\le k}i^d$ 在 $0\le d\le p-2$ 时取 $0,\ldots,d+1$ 共 $d+2$ 个样本。

\section{阶乘表：组合数与排列数}\label{knowledge-factorial-table}\index{阶乘表}\index{排列数}
素数 $p$，$0\le k\le n<p$：
\[
C(n,k)=n!\,(k!)^{-1}\,((n-k)!)^{-1},\qquad
P(n,k)=n!\,((n-k)!)^{-1}\pmod p.
\]
扩表到 $N$：正向补阶乘，只对 $N!$ 求一次逆，再 $(i-1)!^{-1}=i\,(i!)^{-1}$ 反向补。
查询前显式扩表，查询 $O(1)$；扩表后旧元素引用失效。见第~\pageref{compact-Binomial}~页。
